\documentclass[a4paper, serif, 10pt]{moderncv}

%% moderncv
% color and style
\moderncvcolor{black}
\moderncvstyle{banking}

%% page layout/margins
\usepackage[top=2.5cm, bottom=2.5cm, left=1.5cm, right=1.5cm]{geometry}

\nopagenumbers{}

%% Babel config for Norwegian language
% ref:
% - https://latex3.github.io/babel/guides/locale-norwegian.html
\usepackage[norsk]{babel}

%% fontspec
\usepackage{fontspec}

\IfFontExistsTF{Linux Libertine}{
  \setmainfont[Ligatures={TeX, Common}, Numbers=OldStyle]{Linux Libertine}
}{}
\IfFontExistsTF{Linux Libertine O}{
  \setmainfont[Ligatures={TeX, Common}, Numbers=OldStyle]{Linux Libertine O}
}{}

%% CTeX
% CJK support, used to show CJK characters in the resume
%
% - heading=false: disable ctex heading
% - punct=kaiming: use kaiming punctuations styles for CJK
% - scheme=plain: use plain scheme, do not override \normalsize font size
% - space=auto: space settings for CJK characters
%
% CJK fonts are resolved per locale language, e.g. Noto Serif/Sans CJK JP
% for Japanese resumes, with fallback to the next region if not installed
%
% ref:
% - http://ctan.mirrorcatalogs.com/language/chinese/ctex/ctex.pdf
\usepackage[UTF8, heading=false, punct=kaiming, scheme=plain, space=auto]{ctex}

\IfFontExistsTF{Noto Serif CJK SC}{
  \setCJKmainfont{Noto Serif CJK SC}
}{}
\IfFontExistsTF{Noto Sans CJK SC}{
  \setCJKsansfont{Noto Sans CJK SC}
}{}

%% line spacing
\usepackage{setspace}
\setstretch{1.125}

%% URL styling - use normal text instead of monospace
\urlstyle{same}

%% Patch \httplink and \httpslink to handle full URLs
% The original moderncv macros always prepend http:// or https:// to URLs.
% This causes issues when the URL already has a protocol (e.g., https://example.com)
% resulting in malformed URLs like https://https://example.com.
% This patch checks if the URL already contains "://" and uses it directly if so.
% Note: We use \str_set:Nx (x-expansion) to expand macros like \@homepage first.
\ExplSyntaxOn
\str_new:N \g_yamlresume_url_str

\RenewDocumentCommand{\httpslink}{O{}m}{
  \str_set:Nx \g_yamlresume_url_str {#2}
  \str_if_in:NnTF \g_yamlresume_url_str {://}
    { \url{#2} }
    { \url{https://#2} }
}

\RenewDocumentCommand{\httplink}{O{}m}{
  \str_set:Nx \g_yamlresume_url_str {#2}
  \str_if_in:NnTF \g_yamlresume_url_str {://}
    { \url{#2} }
    { \url{http://#2} }
}
\ExplSyntaxOff

\name{Ingrid Solberg}{}
\title{QA-ingeniør med fokus på testautomatisering og kvalitetssikring}
\phone[mobile]{+47 912 34 567}
\email{ingrid.solberg@example.no}
\homepage{https://www.ingridsolberg.no}

\address{Storgata 42, Oslo, Oslo, Norge, 0182}{}{}

\extrainfo{{\small \faGithub}\ \href{https://github.com/ingridsolbergqa}{@ingridsolbergqa} {} {} {} • {} {} {} 
{\small \faLinkedin}\ \href{https://www.linkedin.com/in/ingrid-solberg-qa}{@ingrid-solberg-qa}}

\begin{document}

\maketitle

\section{Grunnleggende informasjon}

\cvline{}{\begin{itemize}
\item Erfaren QA-ingeniør med seks års bakgrunn fra test av web- og API-baserte helseløsninger
\item Sterk på testautomatisering med Playwright, Cypress og Postman, og på å bygge robuste regresjonssuiter
\item Systematisk og kvalitetsbevisst, med erfaring fra Scrum, CI/CD og tverrfaglig produktutvikling
\item Opptatt av tydelig teststrategi, gode feilrapporter og kontinuerlig forbedring av kvalitetsprosesser
\end{itemize}}
\section{Utdanning}

\cventry{aug. 2015–juni 2018}
        {Bachelor, Informatikk: design, bruk og interaksjon, Poeng: 4.2}
        {Universitetet i Oslo}
        {\url{https://www.uio.no}}
        {}
        {\begin{itemize}
\item Bygget solid grunnlag i systemutvikling, databaser og brukersentrert design
\item Gjennomførte testprosjekter med både manuell og automatisert testing
\item Arbeidet i tverrfaglige studentteam med kravanalyse, utvikling og kvalitetssikring
\end{itemize}
\textbf{Kurs}: Programmering og systemutvikling, Algoritmer og datastrukturer, Databaser, Testing av programvare, Webutvikling, Prosjektstyring}
\section{Arbeidserfaring}

\cventry{mars 2021–Nå}
        {QA-ingeniør}
        {Nordisk HelseTech AS}
        {\url{https://www.nordiskhelsetech.no}}
        {}
        {\begin{itemize}
\item Ansvarlig for teststrategi og kvalitetssikring av pasientrettet webplattform og tilhørende API-er
\item Bygget og vedlikeholdt automatiserte regresjonssuiter i Playwright og Cypress, med daglig kjøring i CI/CD
\item Utførte API-testing med Postman og REST-assured, inkludert validering av autentisering, datakvalitet og feilhåndtering
\item Ledet testgjennomganger og triage av feil i tett samarbeid med utviklere, produktleder og designere
\item Forbedret feilrapporteringsmaler og testdokumentasjon, noe som reduserte gjennomsnittlig tid til retting
\end{itemize}
\textbf{Nøkkelord}: Teststrategi, Testautomatisering, Playwright, Cypress, API-testing, CI/CD, Regresjonstesting, Kvalitetssikring}

\cventry{aug. 2018–feb. 2021}
        {Testautomatiseringsutvikler}
        {Oslo Digital AS}
        {\url{https://www.oslodigital.no}}
        {}
        {\begin{itemize}
\item Utviklet automatiserte tester for nettbaserte tjenester innen offentlig sektor
\item Skrev testkode i JavaScript og Python, og integrerte testene i Azure DevOps-pipelines
\item Gjennomførte manuell utforskende testing og sikret god dekning av kritiske brukerreiser
\item Bidro til å innføre BDD med Cucumber og tydeligere akseptansekriterier i Scrum-teamet
\item Dokumenterte testresultater og presenterte kvalitetsstatus i sprintgjennomganger
\end{itemize}
\textbf{Nøkkelord}: JavaScript, Python, Azure DevOps, BDD, Cucumber, Utforskende testing, Scrum}
\section{Språk}

\cvline{Norsk}{Morsmål eller tospråklig nivå \hfill \textbf{Nøkkelord}: Morsmål, Fagtekster}
\cvline{Engelsk}{Fullt profesjonelt nivå \hfill \textbf{Nøkkelord}: TOEFL 105, Teknisk dokumentasjon}
\section{Ferdigheter}

\cvline{Testing og kvalitet}{Ekspert \hfill \textbf{Nøkkelord}: Teststrategi, Testautomatisering, Regresjonstesting, API-testing, Ytelsestesting, Utforskende testing}
\cvline{Verktøy}{Avansert \hfill \textbf{Nøkkelord}: Playwright, Cypress, Postman, JMeter, Git, Jira, Azure DevOps}
\cvline{Programmering}{Viderekommen \hfill \textbf{Nøkkelord}: JavaScript, TypeScript, Python, SQL, Bash}
\cvline{Metodikk}{Avansert \hfill \textbf{Nøkkelord}: Scrum, Agile, CI/CD, TDD, BDD, Risikobasert testing}
\section{Utmerkelser}

\cventry{nov. 2023}
        {Nordisk HelseTech AS}
        {Årets kvalitetspris}
        {}
        {}
        {Tildelt for arbeidet med å redusere produksjonsfeil og øke testdekningen på pasientplattformen gjennom automatisering og tydeligere kvalitetsrutiner.}
\section{Sertifikater}

\cventry{mai 2019}
        {ISTQB}
        {ISTQB-sertifisert tester, grunnleggende nivå}
        {\url{https://www.istqb.org/certifications/certified-tester-foundation-level}}
        {}
        {}

\cventry{sep. 2020}
        {ISTQB}
        {ISTQB Agile Tester Extension}
        {\url{https://www.istqb.org/certifications/agile-tester-extension}}
        {}
        {}
\section{Prosjekter}

\cventry{jan. 2022–des. 2023}
        {Regresjonssuite som verifiserer kritiske pasientreiser og API-kontrakter ved hver utrulling}
        {Automatisert regresjonssuite for pasientplattform}
        {\url{https://www.nordiskhelsetech.no}}
        {}
        {\begin{itemize}
\item Designet testpyramide og prioriterte kritiske brukerreiser for automatisering
\item Reduserte manuell regresjonstid fra tre dager til under fire timer
\item Etablerte rapportering av testresultater og sporbarhet mot krav i Jira
\end{itemize}
\textbf{Nøkkelord}: Playwright, API-testing, Regresjonstesting, CI/CD}

\cventry{apr. 2020–okt. 2020}
        {Dashboard som samler testresultater, feiltrender og dekningsgrad for produktteam}
        {Kvalitetsdashboard for tverrfaglige team}
        {}
        {}
        {\begin{itemize}
\item Bygget en enkel løsning for å visualisere teststatus fra Azure DevOps
\item Ga teamene raskere innsikt i flaky tester og gjentakende feil
\item Forbedret beslutningsgrunnlaget i sprintplanlegging og release-møter
\end{itemize}
\textbf{Nøkkelord}: Azure DevOps, JavaScript, Testmetrikk, Rapportering}
\section{Interesser}

\cvline{Friluftsliv}{Fjellturer, Langrenn, Kajakk}
\cvline{Teknologi}{Åpen kildekode, Testmiljøer, Podkaster}
\section{Frivillig arbeid}

\cventry{jan. 2022–Nå}
        {Frivillig arrangør}
        {Oslo QA-community}
        {\url{https://www.meetup.com/oslo-qa-community/}}
        {}
        {\begin{itemize}
\item Arrangerer månedlige fagmøter om testautomatisering, kvalitet og sikkerhet
\item Koordinerer foredragsholdere og fasiliterer erfaringsutveksling mellom QA-miljøer i Oslo
\item Bidrar til et inkluderende fagmiljø for juniorer og erfarne testere
\end{itemize}}
    
\end{document}